% ================================================================
% RESEARCH MEMO: GEOMETRIC CHARACTERIZATIONS OF CHIRAL KOSZULNESS
% Agent: BLUE-2 (Derived Algebraic Geometer)
% Date: 2026-03-17
% ================================================================

\documentclass[11pt]{article}
\usepackage{amsmath,amssymb,amsthm,mathrsfs}
\usepackage[margin=1.2in]{geometry}
\usepackage{enumitem}

\newcommand{\cA}{\mathcal{A}}
\newcommand{\cC}{\mathcal{C}}
\newcommand{\cD}{\mathcal{D}}
\newcommand{\cO}{\mathcal{O}}
\newcommand{\cZ}{\mathcal{Z}}
\newcommand{\bD}{\mathbb{D}}
\newcommand{\bG}{\mathbb{G}}
\newcommand{\bP}{\mathbb{P}}
\newcommand{\bQ}{\mathbb{Q}}
\newcommand{\bR}{\mathbb{R}}
\newcommand{\bC}{\mathbb{C}}
\newcommand{\fg}{\mathfrak{g}}
\newcommand{\Ran}{\mathrm{Ran}}
\newcommand{\FM}{\mathrm{FM}}
\newcommand{\Conf}{\mathrm{Conf}}
\newcommand{\Trop}{\mathrm{Trop}}
\newcommand{\barB}{\bar{B}}
\newcommand{\MC}{\mathrm{MC}}
\newcommand{\Hom}{\mathrm{Hom}}
\newcommand{\Ext}{\mathrm{Ext}}
\newcommand{\gr}{\mathrm{gr}}
\newcommand{\Sym}{\mathrm{Sym}}
\newcommand{\Def}{\mathrm{Def}}
\newcommand{\Fact}{\mathrm{Fact}}
\newcommand{\Ainf}{A_\infty}
\newcommand{\Linf}{L_\infty}
\newcommand{\Einf}{E_\infty}
\newcommand{\Etwo}{E_2}

\theoremstyle{plain}
\newtheorem{conjecture}{Conjecture}[section]
\newtheorem{proposition}[conjecture]{Proposition}
\newtheorem{question}[conjecture]{Question}
\theoremstyle{definition}
\newtheorem{definition}[conjecture]{Definition}
\theoremstyle{remark}
\newtheorem{remark}[conjecture]{Remark}

\title{Geometric Characterizations of Chiral Koszulness:\\
  Six Conjectures and a Structural Question}
\author{BLUE-2 Research Memo}
\date{March 2026}

\begin{document}
\maketitle

\tableofcontents

% ================================================================
\section{Executive summary}
% ================================================================

The monograph characterizes chiral Koszulness \emph{reductively}:
a chiral algebra~$\cA$ is Koszul if its PBW-associated graded
$\gr_F\cA$ is classically Koszul (Theorem~8.3, the PBW criterion).
This is a powerful recognition tool---it covers Heisenberg,
Kac--Moody, Virasoro, $\mathcal{W}_N$, and Yangians---but it
is algebraic in character.  The geometry of the curve~$X$ enters
only through the propagator on~$\overline{C}_2(X)$, not through
the Koszulness condition itself.

This memo proposes six geometric reformulations of chiral
Koszulness, each exploiting a different piece of the monograph's
machinery, and addresses the structural question of whether
Koszulness depends on~$X$.


% ================================================================
\section{Conjecture A: Ran-space formality}
\label{sec:ran-formality}
% ================================================================

\subsection{Statement}

\begin{conjecture}[Ran-space formality criterion]
\label{conj:ran-formality}
Let $\cA$ be an augmented chiral algebra on a smooth curve~$X$.
The following are equivalent:
\begin{enumerate}[label=(\roman*)]
\item $\cA$ is chiral Koszul (Definition~8.5, the monograph's
  def:chiral-koszul-morphism).
\item The factorization algebra $\cA|_{\Ran^{\leq n}(X)}$ is
  \emph{formal} as an object of
  $D^b(\Fact^{\mathrm{aug}}(X))$ for all~$n$: i.e., it is
  quasi-isomorphic to its cohomology equipped with the
  transferred $\Ainf$-structure, and the transferred operations
  $\tilde{m}_k = 0$ for $k \geq 3$.
\item The bar spectral sequence
  $E_1^{p,q} \Rightarrow H^{p+q}(\Omega(\barB(\cA)))$ collapses
  at $E_2$ (equivalently, the cobar-bar is formal with vanishing
  higher Massey products).
\end{enumerate}
\end{conjecture}

\subsection{Evidence from the monograph}

The equivalence (i)$\Leftrightarrow$(iii) is essentially proved.
Proposition~\texttt{prop:e2-collapse-formality}
(higher\_genus\_complementarity.tex, line~4010) establishes the
equivalence of $E_2$-collapse, formality of
$\Omega_g(\barB_g(\cA))$, and vanishing of transferred operations
$\tilde{m}_n^{(g)}$ for $n \geq 3$.  That proposition operates at
each genus~$g$; the genus-$0$ case is exactly the PBW criterion
(Theorem~\texttt{thm:pbw-koszulness-criterion}).

The new content of Conjecture~\ref{conj:ran-formality} is the
\emph{stratum-by-stratum} interpretation~(ii): formality should
hold not just globally on $\Ran(X)$ but on each finite truncation
$\Ran^{\leq n}(X)$, reflecting the factorization structure.

\subsection{Relation to $\Etwo$-formality}

Kontsevich's formality theorem (Theorem~\texttt{thm:e2-formality},
en\_koszul\_duality.tex) says the $\Etwo$-operad is formal.
Chiral algebras on a curve carry an $\Etwo$-structure from the
little-disks operad on $\bC \subset X$.  The PBW criterion reduces
chiral Koszulness to formality of this $\Etwo$-structure on the
associated graded.  In the geometric picture:
\begin{quote}
\emph{Chiral Koszulness = the factorization algebra inherits
$\Etwo$-formality from its PBW-associated graded.}
\end{quote}

\subsection{Obstacles}

The passage from stratum-level formality to global Ran-space
formality requires a descent argument: formality on each
$\Ran^{\leq n}$ must be compatible with the factorization maps.
This is a nontrivial condition---it asks that the formality
quasi-isomorphisms respect the diagonal maps
$\Ran^{\leq n} \to \Ran^{\leq m}$ for $m < n$.  Such
coherent formality is available for $\Etwo$-algebras
(Tamarkin--Kontsevich), but the chiral extension to
factorization algebras on arbitrary curves is not in the
literature.

\subsection{Utility}

If true, this would give a \emph{cohomological} criterion for
Koszulness: instead of checking acyclicity of twisted tensor
products (an algebraic condition), one checks formality of a
constructible sheaf (a geometric condition).  This is more natural
from the BD perspective and may be computable via mixed Hodge theory
on configuration spaces.


% ================================================================
\section{Conjecture B: Boundary acyclicity on FM}
\label{sec:fm-boundary}
% ================================================================

\subsection{Statement}

\begin{conjecture}[FM boundary acyclicity criterion]
\label{conj:fm-boundary}
Let $\cA$ be a chiral algebra on $X$, and let
$\barB_n^{\mathrm{ch}}(\cA)$ be the $n$-th bar component, a
regular holonomic $\cD$-module on $\overline{C}_n(X)$
(Lemma~\texttt{lem:bar-holonomicity}).  Then $\cA$ is chiral
Koszul if and only if for all $n \geq 2$:
\[
H^k\bigl(i_S^!\, \barB_n^{\mathrm{ch}}(\cA)\bigr) = 0
\quad\text{for all } k \neq 0
\quad\text{and all boundary strata } S \subset \partial\overline{C}_n(X),
\]
where $i_S \colon S \hookrightarrow \overline{C}_n(X)$ is the
inclusion.  In words: the restriction of the bar complex to every
boundary stratum of the FM compactification is acyclic outside
degree~$0$.
\end{conjecture}

\subsection{Structural basis}

The FM boundary strata decompose as products of lower FM spaces:
a stratum where $k$ points collide is locally isomorphic to
$\overline{C}_k(\bC) \times \overline{C}_{n-k+1}(X)$.  This
decomposition is the operadic composition.  The bar differential
\[
d_{\mathrm{bar}} = d_{\mathrm{internal}} + d_{\mathrm{residue}}
  + d_{\mathrm{form}}
\]
(Definition~\texttt{def:bar-differential-complete}) has its
residue component $d_{\mathrm{residue}} = \sum_D \mathrm{Res}_D$
supported exactly on boundary divisors.  The claim is that
Koszulness = \emph{each residue contribution is individually
acyclic}, not merely that their sum is.

\subsection{Evidence}

The PBW criterion shows that chiral Koszulness reduces to
classical Koszulness of $\gr_F\cA$, and classical Koszulness is
equivalent to diagonal Ext-vanishing:
$\Ext^{i,j}_{\gr\cA}(\bC,\bC) = 0$ for $i \neq j$.  In the
geometric picture, $\Ext^{i,j}$ is computed by cohomology on FM
boundary strata of codimension~$i$ with coefficient in weight~$j$.
Diagonal vanishing translates to: cohomology of each boundary
stratum is concentrated in the expected degree, i.e., acyclicity
in the sense of the conjecture.

The bar complex components are regular holonomic $\cD$-modules
(Lemma~\texttt{lem:bar-holonomicity}), so the restriction
$i_S^!\,\barB_n$ is a holonomic module on~$S$.  Acyclicity of
holonomic restrictions is a tameness condition that fits naturally
into the BD framework.

\subsection{Obstacles}

The main difficulty is that the FM boundary has many irreducible
components (indexed by nested partitions), and the acyclicity
condition must hold for \emph{all} of them simultaneously.
This is stronger than diagonal Ext-vanishing, which is an
aggregate statement.  One needs to show that the spectral
sequence from the boundary filtration degenerates, which is
plausible but requires either a purity argument or explicit
computation stratum by stratum.

\subsection{Utility}

A boundary-stratum criterion would make Koszulness visible in
intersection cohomology of the FM space.  It would also connect
directly to the weight filtration on the convolution algebra
$\fg^{\mathrm{mod}}_\cA$, whose strata \emph{are} the FM boundary
strata (Definition~\texttt{def:weight-filtration-tower}).


% ================================================================
\section{Conjecture C: Tropical Koszulness via planted forests}
\label{sec:tropical}
% ================================================================

\subsection{Statement}

\begin{conjecture}[Tropical characterization of Koszulness]
\label{conj:tropical-koszul}
Let $(X,D)$ be a smooth curve with snc divisor, and let
$\bG_{\mathrm{pf}}$ be the planted-forest coefficient algebra
(Definition~\texttt{def:planted-forest-coefficient-algebra}).
Define the \emph{tropical bar complex} of $\cA$ by
\[
\barB^{\mathrm{trop}}_n(\cA)
:= \bigoplus_{F \in \mathsf{For}_{\mathrm{plant}}(n)}
   \kappa_F(\cA) \cdot e_F
\;\in\;
\Def_{\mathrm{cyc}}(\cA) \;\widehat{\otimes}\; \bG_{\mathrm{pf}},
\]
where $\kappa_F(\cA)$ is the weight of the shadow tower
$\Theta^{\leq r}_\cA$ on the planted forest~$F$.
Then $\cA$ is chiral Koszul if and only if
$\barB^{\mathrm{trop}}(\cA)$ is acyclic for the differential
$d_{\mathrm{pf}}$ (edge contraction, the planted-forest differential) tensored with the
deformation differential on $\Def_{\mathrm{cyc}}(\cA)$.
\end{conjecture}

\subsection{Structural basis}

By Mok's theorem (Theorem~\texttt{thm:planted-forest-tropicalization}),
the planted-forest coefficient algebra $\bG_{\mathrm{pf}}$ is the
chain algebra on the face poset of
$\Trop(\FM_n(X|D))$.  The codimension formula
\[
\mathrm{codim}(F) = \sum w_i + \sum(|V(T_{v_j})| - 1)
\]
(Proposition~\texttt{prop:planted-forest-tropical}, equation~(8932))
matches the weight filtration on $\fg^{\mathrm{amb}}_\cA$.
The shadow Postnikov tower $\Theta^{\leq r}_\cA$ already encodes
the contributions of planted forests of bounded depth~$r$.

\subsection{The tropical Koszul complex}

At the tropical level, the bar differential becomes purely
combinatorial: it is edge contraction in the planted-forest
category, with coefficients in the cyclic deformation complex.
The Koszul condition would then be:
\begin{quote}
\emph{The tropical bar complex, graded by forest depth, has
cohomology concentrated on the diagonal (depth = cohomological
degree).}
\end{quote}
This is an intrinsically combinatorial condition on the weights
$\kappa_F(\cA)$ assigned to planted forests by the shadow tower.

\subsection{Evidence}

For one-channel families (Heisenberg, Kac--Moody), the shadow
tower satisfies scalar saturation:
$\Theta_\cA = \kappa \cdot \eta \otimes \Lambda$.  In this case,
the tropical bar complex collapses because all higher-order forest
contributions are determined by the single scalar~$\kappa$, and
the resulting complex is exact (it reduces to the classical Koszul
complex for~$\Sym(V)$).

For the Virasoro algebra, the shadow tower is \emph{infinite}:
the quintic obstruction $o^{(5)}_{\mathrm{Vir}} \neq 0$, forcing
an infinite cascade of corrections.  Yet Virasoro is still Koszul
(Theorem~\texttt{thm:virasoro-chiral-koszul}).  This means the
tropical complex is acyclic despite having contributions at all
forest depths---an interesting constraint on the interplay between
the quartic resonance class $Q^{\mathrm{contact}}_{\mathrm{Vir}}
= 10/[c(5c+22)]$ and the higher obstructions.

\subsection{Obstacles}

The main obstacle is that the tropical bar complex captures only
the \emph{combinatorial shadow} of the full bar complex.  The
passage from the algebraic bar complex on~$\FM_n(X|D)$ to its
tropicalization loses analytic information (e.g., periods,
monodromy).  One needs a ``tropical descent'' theorem: acyclicity
of the tropical complex implies acyclicity of the algebraic one.
This would follow from a tropical analog of the Deligne weight
spectral sequence, but such a tool is not yet available in the
log~FM setting.

\subsection{Utility}

A tropical criterion would make Koszulness effectively computable
for new examples: one would only need to evaluate the shadow tower
on finitely many planted forests (up to a given depth) and check
a combinatorial acyclicity condition.  This could automate the
verification of modular Koszulness at higher genus.


% ================================================================
\section{Conjecture D: D-module characterization}
\label{sec:d-module}
% ================================================================

\subsection{Statement}

\begin{conjecture}[$\cD$-module tameness criterion]
\label{conj:d-module}
Let $\cA$ be a chiral algebra on a smooth projective curve~$X$,
finitely generated over~$\cD_X$ with regular singularities.
Then $\cA$ is chiral Koszul if and only if:
\begin{enumerate}[label=(\roman*)]
\item Each bar component $\barB_n^{\mathrm{ch}}(\cA)$ has
  \emph{pure} weight on $\overline{C}_n(X)$ (in the sense of
  mixed Hodge modules).
\item The characteristic variety
  $\mathrm{Ch}(\barB_n^{\mathrm{ch}}(\cA)) \subset T^*\overline{C}_n(X)$
  is contained in the union of conormal bundles to the boundary
  strata: $\mathrm{Ch}(\barB_n) \subset \bigcup_S T^*_S \overline{C}_n$.
\end{enumerate}
\end{conjecture}

\subsection{Structural basis}

The monograph proves that each bar component is a regular holonomic
$\cD$-module on $\overline{C}_n(X)$
(Lemma~\texttt{lem:bar-holonomicity}).  Holonomicity means the
characteristic variety has dimension $= \dim \overline{C}_n$.
Condition~(ii) above is a refinement: it asks that the singular
support is ``aligned'' with the FM boundary stratification.

This is natural from the BD perspective.  Chiral algebras are
$\cD$-modules on~$X$; the bar construction produces $\cD$-modules
on configuration spaces.  The OPE singularities are encoded in
the characteristic variety.  Koszulness is the condition that
these singularities are ``as simple as possible''---specifically,
that the bar $\cD$-module is a successive extension of rank-one
connections along boundary strata, which is the $\cD$-module
translation of diagonal Ext-vanishing.

\subsection{Purity and Koszulness}

Condition~(i) connects to a deep observation: for classical
Koszul algebras, the bar complex has \emph{pure} Hodge-theoretic
weight.  (This is the content of Backelin's purity theorem:
Koszulness $\Leftrightarrow$ the bar complex is pure as a mixed
Hodge structure.)  Condition~(i) lifts this to the chiral
setting: purity of $\barB_n^{\mathrm{ch}}(\cA)$ as a mixed
Hodge module on~$\overline{C}_n(X)$.

\subsection{Obstacles}

\begin{itemize}
\item Mixed Hodge modules on FM spaces are technically demanding.
  Saito's theory applies, but explicit computations are rare.
\item The purity condition (i) may be too strong: classical
  Koszulness implies purity, but the converse requires
  finite-dimensionality conditions that may fail in the
  chiral setting (where weight spaces are infinite-dimensional
  in general, though finite in each conformal weight).
\item Condition (ii) is a ``microlocal'' criterion.  Verifying
  it requires computing the singular support of the bar
  $\cD$-module, which in practice means understanding the
  wavefront set of the integration kernel on $X^n \setminus
  \Delta$.
\end{itemize}

\subsection{Utility}

A $\cD$-module criterion would connect Koszulness to
representation-theoretic invariants: the characteristic variety
encodes the ``microlocal support'' of the chiral algebra, and
conditions on it translate into finiteness conditions on category
$\cO$.  This could give a path to proving Koszulness for new
families (e.g., exceptional $\mathcal{W}$-algebras) via
characteristic variety computations rather than explicit PBW
reductions.


% ================================================================
\section{Conjecture E: Shifted-symplectic characterization}
\label{sec:shifted-symplectic}
% ================================================================

\subsection{Statement}

\begin{conjecture}[Lagrangian criterion for Koszulness]
\label{conj:lagrangian}
Let $(\cA, \cA^!)$ be a chiral pair.  Suppose the ambient
complementarity tangent complex $T_{\mathrm{comp}}(\cA)$
(Definition~\texttt{def:ambient-tangent-complex}) carries a
$(-1)$-shifted symplectic structure $\omega$ and the one-sided
deformation maps
$i_\cA \colon \mathcal{M}_\cA \to \mathcal{M}_{\mathrm{comp}}$
and
$i_{\cA^!} \colon \mathcal{M}_{\cA^!} \to \mathcal{M}_{\mathrm{comp}}$
are isotropic.  Then:
\begin{enumerate}[label=(\roman*)]
\item $\cA$ is chiral Koszul if and only if $i_\cA$ and
  $i_{\cA^!}$ are \emph{Lagrangian} (isotropic of maximal
  dimension).
\item The Lagrangian condition is equivalent to: the bar-cobar
  counit $\Omega(\barB(\cA)) \to \cA$ is a quasi-isomorphism.
\end{enumerate}
In other words: \emph{Koszulness = the bar-cobar correspondence
is a Lagrangian correspondence in the $(-1)$-shifted symplectic
deformation space.}
\end{conjecture}

\subsection{Evidence from the monograph}

The ambient complementarity theorem
(Conjecture~\texttt{thm:ambient-complementarity} in
chiral\_hochschild\_koszul.tex) already conjectures (as three
sub-conjectures: self-duality, one-sided isotropy, ambient
complementarity) that $T_{\mathrm{comp}}(\cA)$ carries a
$(-1)$-shifted symplectic structure and that the one-sided maps
are Lagrangian.  The evidence sections in the monograph verify
isotropy by differentiating the MC equation.

The proved linear shadow is Theorem~C: $Q_g(\cA) \oplus Q_g(\cA^!)
\cong H^*(\overline{\mathcal{M}}_g, \cZ(\cA))$ as complementary
Lagrangian subspaces
(Theorem~\texttt{thm:quantum-complementarity-main}).  The
Lagrangian upgrade is conditional on perfectness of the cyclic
deformation complex.

\subsection{The PTVV framework}

The $(-1)$-shifted symplectic geometry of bar-cobar duality
fits naturally into the PTVV framework
(Pantev--To\"{e}n--Vaqui\'{e}--Vezzosi).  In that framework,
a $(-1)$-shifted symplectic structure on a derived stack~$\mathcal{M}$
produces Lagrangian intersections that compute derived tensor
products.  The bar-cobar adjunction is an intersection in
$\mathcal{M}_{\mathrm{comp}}$: the derived intersection
$\mathcal{M}_\cA \times^h_{\mathcal{M}_{\mathrm{comp}}}
\mathcal{M}_{\cA^!}$ should compute the twisted tensor product
$K_\tau(\cA, \cA^!)$, and its acyclicity (= Koszulness) is the
condition that the intersection is \emph{transverse}---i.e., that
the Lagrangians are complementary.

\subsection{Obstacles}

\begin{itemize}
\item The ambient complementarity theorem is currently
  \textbf{conjectural} in the monograph (three sub-conjectures).
  The Lagrangian criterion for Koszulness depends on establishing
  it.
\item The passage from $(-1)$-shifted symplectic geometry to
  Koszulness requires understanding the non-linear structure of
  $\mathcal{M}_{\mathrm{comp}}$, not just its tangent complex.
  The PTVV machinery handles derived Artin stacks; the chiral
  deformation space may be more singular.
\item The cyclic pairing on $\Def_{\mathrm{cyc}}(\cA)$ needs to
  be \emph{perfect} for the Lagrangian condition to be meaningful.
  Perfectness is a non-trivial condition that may fail for
  infinite-type algebras ($\mathcal{W}_\infty$).
\end{itemize}

\subsection{Utility}

If Koszulness = Lagrangian complementarity, then the failure of
Koszulness has a clean geometric interpretation: the two
deformation spaces fail to be transverse, which means there is a
common tangent direction---a \emph{simultaneous deformation} of
$\cA$ and $\cA^!$ that is not captured by deforming either one
alone.  This would give a conceptual explanation for why
Koszulness fails off the Koszul locus: the coderived category
intervenes precisely when the Lagrangians fail to be transverse.


% ================================================================
\section{Conjecture F: Deformation quantization compatibility}
\label{sec:dq}
% ================================================================

\subsection{Statement}

\begin{conjecture}[Quantization-compatible Koszulness]
\label{conj:dq}
Let $\cA_\hbar$ be a one-parameter deformation of a commutative
chiral algebra $\cA_0$ (i.e., a vertex Poisson algebra / coisson
algebra).  Assume $\cA_0$ is classically Koszul.  Then $\cA_\hbar$
is chiral Koszul if and only if:
\begin{enumerate}[label=(\roman*)]
\item The deformation is \emph{Koszul-compatible}: the PBW
  filtration on~$\cA_\hbar$ is a flat deformation of the weight
  filtration on~$\cA_0$.
\item The quantum correction $\cA_\hbar - \cA_0$ lies in
  $F_{\geq 1}\cA_\hbar$ (the deformation does not change the
  degree-$0$ part of the filtration).
\end{enumerate}
\end{conjecture}

\subsection{Structural basis}

The PBW criterion (Theorem~\texttt{thm:pbw-koszulness-criterion})
already says: if $\gr_F\cA$ is classically Koszul and the
filtration is flat, then $\cA$ is chiral Koszul.  The conjecture
above reverses this: it asks whether Koszul-compatible quantization
is the \emph{only} way to produce chiral Koszul algebras.

In the monograph's four-regime hierarchy, every chiral Koszul
algebra sits in one of: quadratic ($\mu_0 = 0$), curved-central,
filtered-complete, or programmatic.  The first three all have
PBW filtrations; the fourth ($\mathcal{W}_\infty$) is the MC4
completion target.  The conjecture asks: does every chiral Koszul
algebra admit a PBW filtration whose associated graded is a
classically Koszul vertex Poisson algebra?

\subsection{Relation to the coisson shadow}

In the three-pillar architecture (Pillar~A, MS24), the primitive
local object is a homotopy chiral algebra $\mathrm{Ch}_\infty$,
with strict chiral algebra appearing after rectification, and the
coisson shadow as $\mathrm{Ch}_\infty \to \text{strict on }
H^\bullet \to \text{PVA/coisson}$.  The conjecture asks whether
this chain respects Koszulness:
\[
\cA_\hbar \;\text{chiral Koszul}
\quad\Longleftrightarrow\quad
\cA_0 = \gr_F \cA_\hbar \;\text{classically Koszul}
\quad+\quad
\text{flatness of } F.
\]

\subsection{Obstacles}

The conjecture could fail if there exist chiral Koszul algebras
that do not admit \emph{any} PBW filtration with classically Koszul
associated graded.  This would require a genuinely non-PBW proof
of Koszulness.  No such example is currently known, but the
programmatic regime ($\mathcal{W}_\infty$) is a candidate: the
MC4 completion involves infinitely many generators, and whether
the limit admits a PBW filtration is not obvious.

\subsection{Utility}

If true, this conjecture would reduce chiral Koszulness entirely
to the classical theory: one would need only check classical
Koszulness of the vertex Poisson algebra $\cA_0$ (a well-understood
problem) and flatness of the deformation (a condition on the
Rees construction).  It would also connect chiral Koszulness to
the theory of Rees algebras and filtered derived categories.


% ================================================================
\section{The key geometric question: does $X$ matter?}
\label{sec:curve-dependence}
% ================================================================

\subsection{Genus-0 independence}

\begin{proposition}[Koszulness is genus-0 independent of $X$]
\label{prop:genus-0-independence}
At genus~$0$, chiral Koszulness of $\cA$ on $X$ depends only on
the vertex algebra $V = \Gamma(X, \cA)$, not on the curve~$X$.
\end{proposition}

This is clear from the monograph: the PBW criterion
(Theorem~\texttt{thm:pbw-koszulness-criterion}) checks
classical Koszulness of $\gr_F\cA$, which is determined by the
vertex algebra generators and relations data.  The curve~$X$
enters only through the propagator on $\overline{C}_2(X)$, but
the propagator's residue (which determines the bar differential)
is independent of the global geometry of~$X$---it depends only
on the local OPE data.

More precisely: on $\bP^1$ (or any rational curve), the BD
equivalence identifies chiral algebras with vertex algebras, and
the chiral bar construction restricts to the classical bar
construction on generators-and-relations data
(koszul\_pair\_structure.tex, line~95).

\subsection{Higher-genus dependence: the modular Koszulity question}

\begin{question}[Does modular Koszulness depend on $X$?]
\label{q:modular-curve}
Fix a chiral algebra~$\cA$.  Is the modular Koszulity condition
(diagonal Ext-vanishing $\Ext^{i,j}_{\cA^{(g)}}(\bC,\bC) = 0$
for $i \neq j$ at genus~$g \geq 1$) independent of the genus-$g$
curve on which $\cA$ is placed, or does it depend on the choice of
$\Sigma_g \in \overline{\mathcal{M}}_g$?
\end{question}

The monograph's evidence points both ways:

\textbf{Evidence for independence:}
\begin{itemize}
\item The shadow algebra $\cA^{\mathrm{sh}}$ is a homotopy
  invariant (Theorem~\texttt{thm:shadow-homotopy-invariance}):
  quasi-isomorphic chiral algebras have the same shadow algebra,
  and $\kappa(\cA)$ is a quasi-isomorphism invariant.  This
  suggests that Koszulness data is intrinsic to~$\cA$.
\item The curvature $\kappa \cdot \omega_g$ at genus~$g$ factors
  as a product of $\kappa(\cA)$ (depending only on~$\cA$) and
  $\omega_g \in H^2(\overline{\mathcal{M}}_g)$ (depending only
  on the moduli space, not on the choice of curve in it).
\item For Heisenberg and Kac--Moody at generic level, modular
  Koszulness is verified directly from the explicit bar complex,
  which depends on~$\cA$ and the abstract genus~$g$, not on the
  specific curve.
\end{itemize}

\textbf{Evidence for dependence:}
\begin{itemize}
\item The genus-$g$ bar complex lives on $\overline{C}_n(\Sigma_g)$,
  whose topology depends on $\Sigma_g$.  The propagator on
  $\overline{C}_2(\Sigma_g)$ depends on the conformal structure
  (through the prime form or Szeg\H{o} kernel).
\item The Remark~\texttt{rem:genus-graded-obstruction}
  (higher\_genus\_modular\_koszul.tex, line~173) explicitly
  states: ``The modular Koszulity hypothesis is non-trivial:
  genus-$0$ Koszulity does \emph{not} formally imply the
  analogous vanishing at higher genus, where modular forms and
  period integrals introduce new cohomological contributions.''
\item At genus $g \geq 2$, the bar differential depends on period
  integrals of the curve, and these are not constant on
  $\overline{\mathcal{M}}_g$.
\end{itemize}

\subsection{Resolution: stratified independence}

I conjecture that the correct answer involves a \emph{stratified}
independence:

\begin{conjecture}[Stratified curve-independence of Koszulness]
\label{conj:stratified}
\begin{enumerate}[label=(\roman*)]
\item Chiral Koszulness at genus~$0$ is independent of~$X$.
\item Modular Koszulness at genus~$g$ depends on the
  \emph{stratum} of $\overline{\mathcal{M}}_g$ containing the
  curve, but not on the specific curve within a stratum.
  Specifically, on the open stratum $\mathcal{M}_g$ (smooth
  curves), Koszulness is constant; on boundary strata (nodal
  curves), Koszulness may change via the clutching morphism.
\item The change in Koszulness across strata is controlled by the
  curvature: a Koszul pair remains Koszul under degeneration if
  and only if the curvature class $\kappa \cdot \omega_g$
  deforms flatly across the boundary.
\end{enumerate}
\end{conjecture}

\emph{Evidence}: The modular Koszul object
(Definition~\texttt{def:modular-koszul-homotopy}) is defined
\emph{over} $\overline{\mathcal{M}}_g$, and the family statement
in Theorem~A (line~1040--1042 of chiral\_koszul\_pairs.tex) says
the bar-cobar equivalences are ``functorial in families over the
modular configuration spaces $\overline{\mathcal{M}}_{g,n}$.''
This functoriality means Koszulness is a \emph{local} condition
on the base $\overline{\mathcal{M}}_g$: it holds on an open set,
and the question is whether this open set is all
of~$\overline{\mathcal{M}}_g$ or only part of it.

The deepest evidence comes from the center local system
$\cZ_\cA$ on $\overline{\mathcal{M}}_{g,\bullet}$ (item D4
in the modular Koszul object definition): the center \emph{is}
a local system, so its properties are constant on connected
components of~$\overline{\mathcal{M}}_g$.


% ================================================================
\section{Summary table}
\label{sec:summary}
% ================================================================

\begin{center}
\renewcommand{\arraystretch}{1.4}
\begin{tabular}{|c|p{3.8cm}|p{3.8cm}|p{3.8cm}|}
\hline
\textbf{Conj.} & \textbf{Koszulness =} & \textbf{Key obstacle} & \textbf{Utility} \\
\hline
A & Ran-space formality & Coherent descent for formality & Cohomological criterion \\
\hline
B & FM boundary acyclicity & Stratum-by-stratum vs.\ aggregate & Intersection-cohomological \\
\hline
C & Tropical bar acyclicity & Tropical descent theorem & Computational/algorithmic \\
\hline
D & $\cD$-module purity + characteristic variety containment & Mixed Hodge modules on FM & Representation-theoretic \\
\hline
E & Lagrangian complementarity in shifted-symplectic space & Ambient complementarity conjectural & Conceptual (why Koszulness fails) \\
\hline
F & Koszul-compatible quantization of classical Koszul PVA & Existence of PBW filtration for all examples & Reduces to classical theory \\
\hline
\end{tabular}
\end{center}

\medskip
\noindent\textbf{Curve dependence}: genus-$0$ Koszulness is
independent of~$X$; modular Koszulness at genus $g \geq 1$ is
conjecturally constant on open strata of $\overline{\mathcal{M}}_g$
but may change across boundary strata.


% ================================================================
\section{Recommended priorities}
\label{sec:priorities}
% ================================================================

\begin{enumerate}
\item \textbf{Conjecture B (FM boundary acyclicity)} is the most
  immediately tractable: it requires only a careful analysis of the
  spectral sequence from the boundary filtration on the holonomic
  bar complex, and the monograph already has all the ingredients
  (Lemma~\texttt{lem:bar-holonomicity}, the PBW criterion, the
  Verdier extension exchange).

\item \textbf{Conjecture C (tropical Koszulness)} is the most
  novel and has the highest computational payoff.  It requires
  developing the ``tropical descent'' tool in the log~FM setting
  of Mok25, which is a natural next step given the planted-forest
  tropicalization theorem already in the monograph.

\item \textbf{Conjecture E (Lagrangian criterion)} is the deepest
  but depends on resolving the three sub-conjectures of ambient
  complementarity (Conjectures~\texttt{lem:ambient-self-duality},
  \texttt{lem:one-sided-isotropy},
  \texttt{thm:ambient-complementarity}).  If those are proved, the
  Lagrangian criterion follows by PTVV machinery.

\item \textbf{Question~\ref{q:modular-curve} (curve dependence)}
  should be resolved in parallel: a definitive answer would either
  simplify all six conjectures (if Koszulness is always
  curve-independent) or reveal new phenomena (if genus-$g$
  Koszulness can fail on specific subloci of
  $\overline{\mathcal{M}}_g$).
\end{enumerate}

\end{document}
